%% ============================================================
%% Episode 12 - Farsi Podcast Script (Translation)
%% Source: specialized_advisor_report_complete/markdown/12_advisor_defense_qa.md
%% Series: Advisor Progress Report - Deep Dive
%% Compiler: XeLaTeX ONLY
%% Font: B Nazanin
%% CRITICAL: xepersian must always be loaded LAST
%% ============================================================

\documentclass[12pt, a4paper]{article}

%% ============================================================
%% COMPATIBILITY SHIMS
%% ============================================================
\makeatletter
\providecommand\IfClassLoadedT[2]{\@ifclassloaded{#1}{#2}{}}
\providecommand\IfClassLoadedF[2]{\@ifclassloaded{#1}{}{#2}}
\providecommand\IfPackageLoadedT[2]{\@ifpackageloaded{#1}{#2}{}}
\providecommand\IfPackageLoadedF[2]{\@ifpackageloaded{#1}{}{#2}}
\providecommand\IfFileLoadedTF[3]{\@ifpackageloaded{#1}{#2}{#3}}
\providecommand\IfFileLoadedT[2]{\@ifpackageloaded{#1}{#2}{}}
\providecommand\IfFileLoadedF[2]{\@ifpackageloaded{#1}{}{#2}}
\makeatother
\usepackage{expl3}
\ExplSyntaxOn
\cs_if_exist:NF \prop_gput_if_not_in:Nnn
  {
    \cs_new:Npn \prop_gput_if_not_in:Nnn #1#2#3
      { \prop_if_in:NnF #1 {#2} { \prop_gput:Nnn #1 {#2} {#3} } }
  }
\ExplSyntaxOff
%% ============================================================

\usepackage[
  a4paper,
  top=2.5cm, bottom=2.5cm,
  left=2.8cm, right=2.8cm,
  headheight=25pt
]{geometry}

\usepackage{xcolor}
\definecolor{seriesblue}{RGB}{25, 80, 150}
\definecolor{audionote}{RGB}{130, 50, 15}
\definecolor{sectiongray}{RGB}{75, 75, 75}
\definecolor{audiotan}{RGB}{252, 243, 233}

\usepackage{amsmath}
\usepackage{amssymb}
\usepackage{enumitem}
\usepackage{setspace}
\setstretch{1.8}
\usepackage{mdframed}
\usepackage{titlesec}

%% ============================================================
%% xepersian MUST be loaded LAST
%% ============================================================
\usepackage{xepersian}
\settextfont[Scale=1.05]{B Nazanin}
\setlatintextfont{Times New Roman}

\newcommand{\dash}{\lr{---}}

%% ============================================================
%% Page style -- savebox approach
%% ============================================================
\newsavebox{\hdrLeft}
\newsavebox{\hdrRight}
\AtBeginDocument{%
  \savebox{\hdrLeft}{\normalfont\small\color{sectiongray}%
    گزارش پيشرفت مشاور / بررسی عمیق}%
  \savebox{\hdrRight}{\normalfont\small\color{sectiongray}%
    قسمت دوازدهم: تمرين دفاع}%
}
\makeatletter
\def\ps@podcast{%
  \def\@oddhead{\usebox{\hdrLeft}\hfill\usebox{\hdrRight}}%
  \let\@evenhead\@oddhead
  \def\@oddfoot{\normalfont\hfil\thepage\hfil\relax}%
  \let\@evenfoot\@oddfoot
  \let\@mkboth\markboth
}
\makeatother
\pagestyle{podcast}

%% ============================================================
%% Section formatting
%% ============================================================
\titleformat{\section}
  {\large\bfseries\color{seriesblue}}
  {}
  {0pt}
  {}
  [\vspace{-2pt}{\color{seriesblue}\rule{\linewidth}{0.7pt}}\vspace{2pt}]

\titleformat{\subsection}
  {\normalsize\bfseries\color{sectiongray}}
  {}
  {0pt}
  {}

\setcounter{secnumdepth}{0}

\setlist[itemize]{
  itemsep=4pt,
  topsep=5pt,
  parsep=0pt,
  label=$\bullet$
}

%% ============================================================
\begin{document}
%% ============================================================

\thispagestyle{empty}

\vspace*{1.2cm}

{\centering
{\color{seriesblue}\fontsize{24}{30}\selectfont\bfseries قسمت دوازدهم}

\vspace{0.5cm}
{\fontsize{15}{20}\selectfont\bfseries تمرين دفاع از پايان‌نامه: هفت سوال احتمالی}

\vspace{0.9cm}
{\color{seriesblue}\rule{0.55\linewidth}{1pt}}
\vspace{0.7cm}

{\normalsize
\begin{tabular}{rl}
  \textbf{مجموعه:}     & گزارش پيشرفت مشاور \dash{} بررسی عمیق \\[5pt]
  \textbf{مدت:}        & ۱۰ تا ۱۲ دقیقه \\[5pt]
  \textbf{راوی:}       & مجری تنها \\[5pt]
  \textbf{منابع:}      & گزارش کامل مشاور \dash{} بخش‌های ۱ تا ۷ \\
\end{tabular}
}

\vspace{0.9cm}
{\color{seriesblue}\rule{0.55\linewidth}{0.4pt}}
\par}

\vspace{1.2cm}

%% ============================================================
\section{افتتاحیه: سوالات سخت را بشناسيد}

دوازده قسمت کل محتوای فنی گزارش پيشرفت مشاور را پوشش دادند. اين قسمت آخر آن دانش را در يک قالب متفاوت به کار می‌گيرد: سوالات مشاور شبيه‌سازی‌شده و پاسخ‌های گفتاری کامل.

جلسه مشاور يک امتحان با متن ثابت نيست. اما سوالات خاصی تقريباً تضمين‌شده هستند با توجه به محتوای گزارش. نقض بهره استحکام \lr{STA}. اصلاح ۳.۷ درصد. شکست هيبريد. اينها جاهايی هستند که گزارش درباره محدوديت‌هايش صادق‌ترين است \dash{} و صداقت درباره محدوديت‌ها دقيقاً همان چيزی است که يک مشاور دنبال می‌کند و بررسی می‌کند.

هدف اين قسمت ارائه متن برای حفظ کردن نيست. هدف اطمينان از روشن بودن منطق پشت هر پاسخ است تا بتوانيد آن را به طبيعی‌ترين شکل در کلمات خودتان بيان کنيد، حتی اگر سوال از زاويه‌ای غيرمنتظره بيايد.

%% ============================================================
\section{سوال اول: چرا مدل ساده برای بهينه‌سازی و مدل کامل برای معيار؟}

اين احتمالاً اولين سوال در هر بحث فنی است. به نظر مثل يک ناسازگاری می‌رسد.

\textbf{پاسخ:} \lr{PSO} يک الگوريتم جستجو است که تابع برازندگی را ده‌ها هزار بار ارزيابی می‌کند \dash{} ۴۰ ذره ضربدر ۲۰۰ تکرار ۸۰۰۰ ارزيابی به ازای هر کنترل‌گر است. روی مدل کامل غيرخطی، هر شبيه‌سازی ۱۰ ثانيه‌ای ۳ تا ۴ ثانيه طول می‌کشد. يک اجرای \lr{PSO} را به بيش از ۸ ساعت تبديل می‌کند. چهار کنترل‌گر بيش از ۳۰ ساعت طول می‌کشيد.

مدل ساده‌شده \dash{} که معادلات را حول تعادل عمودی خطی می‌کند \dash{} ۱۰ تا ۵۰ برابر سريع‌تر اجرا می‌شود. \lr{PSO} در دقيقه‌ها تمام می‌شود. اين توسعه تکراری و تنظيم مجدد را امکان‌پذير می‌کند.

اما استفاده از مدل ساده‌شده برای نتايج نهايی يعنی معيارهای ما سيستم واقعی ادعاشده را اندازه‌گيری نمی‌کنند. پس \lr{PSO} از مدل ساده برای پيدا کردن بهره‌های خوب به سرعت استفاده می‌کند، و سپس آن بهره‌ها روی مدل کامل غيرخطی برای همه نتايج گزارش‌شده تأييد می‌شوند. هر عدد در جداول عملکرد، هر نتيجه مونت‌کارلو، هر اندازه‌گيری انرژی \dash{} همه از مدل کامل. اين رويکرد استانداردی در بهينه‌سازی مبتنی بر شبيه‌سازی است.

%% ============================================================
\section{سوال دوم: سطح لغزشی خطی را چرا انتخاب کرديد، نه انتگرالی يا پايانی؟}

اين سوال آزمايش می‌کند آيا انتخاب سطح عمدی بود يا پيش‌فرض.

\textbf{پاسخ:} هر دو گزينه بررسی و با دلايل مستند رد شدند. سطح لغزشی انتگرالی يک متغير حالت هفتم به آناليز لياپانوف اضافه می‌کند. تابع لياپانوف مرکب بايد حالت انتگرالی را شامل شود، و اثبات همگرايی به شکل قابل توجهی پيچيده‌تر می‌شود. برای سيستمی که از قبل شش‌بعدی است، پيچيدگی اضافه‌شده بدون فايده متناسب نيست \dash{} خطای ماندگار در شبيه‌سازی ما از قبل کوچک است.

سطح لغزشی پايانی برای همگرايی زمان محدود از يک توان کسری استفاده می‌کند. مشکل يک تکينگی مشتق است: وقتی زاويه آونگ دقيقاً صفر است، مشتق آن توان کسری منفجر می‌شود. روی سخت‌افزار، کوانتيزاسيون حسگر يعنی زاويه هرگز دقيقاً صفر نيست اما در ناحيه مرده دور صفر نوسان می‌کند \dash{} که تکينگی را مکرراً فعال می‌کند.

سطح خطی درجه نسبی دقيقاً يک را می‌دهد، اثبات لياپانوف قابل مديريت است، \lr{PSO} می‌تواند همه چهار پارامتر سطح را همزمان تنظيم کند، و هيچ مشکل عددی در تعادل وجود ندارد. اين توازن مناسب برای يک سيستم در سطح پايان‌نامه است.

%% ============================================================
\section{سوال سوم: بهره‌های استحکام \lr{STA} شرط نظری را نقض نمی‌کنند؟}

اين حساس‌ترين سوال فنی در گزارش است. آن را کوچک نشان ندهيد.

\textbf{پاسخ:} برای بهره‌های اسمی \lr{STA} \dash{} \lr{K-one} مساوی ۸ و \lr{K-two} مساوی ۴ \dash{} شرايط مورنو-اوسوريو کاملاً ارضا شده‌اند. \lr{K-one} از \lr{K-two} بيشتر است، اثبات پايداری اعمال می‌شود، و نتيجه همگرايی زمان محدود معتبر است. همه نتايج معيار منتشرشده برای \lr{STA} از اين بهره‌های اسمی استفاده می‌کنند.

\lr{PSO} تنظيم استحکام \dash{} که از ۱۵ سناريو برای يافتن بهره‌هايی که در طيف گسترده‌ای از شرايط اوليه خوب عمل می‌کنند استفاده می‌کند \dash{} بهره‌هايی با \lr{K-one} مساوی ۲.۰۲ و \lr{K-two} مساوی ۶.۶۷ توليد کرد. اين نقض \lr{K-one} بيشتر از \lr{K-two} است. اين يک مسئله يکپارچگی مستند در گزارش است. بهره‌های استحکام نتايج تجربی هستند \dash{} آن‌ها لرزش را در مجموعه سناريوها کاهش می‌دهند \dash{} اما تضمين پايداری نظری از اثبات فعلی ندارند.

مسير رفع نياز به يکی از دو چيز دارد: اجرای مجدد \lr{PSO} استحکام با قيد \lr{K-one} بيشتر از \lr{K-two} فعال‌شده به صورت محدود در طول جستجو، يا يافتن يک نتيجه پايداری گسترش‌يافته در ادبيات که رژيم \lr{K-one} کمتر از \lr{K-two} را پوشش دهد. نکته کليدی اين است که اين محدوديت فقط بهره‌های استحکام را محدود می‌کند. مشارکت اصلی \lr{STA} \dash{} همگرايی زمان محدود با بهره‌های اسمی، کاهش لرزش ۷۴ درصدی \dash{} تأثير نمی‌گيرد.

%% ============================================================
\section{سوال چهارم: تصديق لياپانوف ۹۶.۲ درصد است نه ۱۰۰ درصد، آيا سيستم واقعاً پايدار است؟}

\textbf{پاسخ:} ۳.۸ درصد نمونه‌هايی که \lr{V-dot} منفی نيست در طول مرحله عبور از لايه مرزی رخ می‌دهند \dash{} به طور خاص وقتی مسير داخل باند جايی است که قدرمطلق سيگما کمتر از اپسيلون است. تابع اشباع دقيقاً برای اجازه دادن به رفتار نرم داخل اين باند به جای سوئيچينگ سخت معرفی شد. داخل لايه مرزی، \lr{V-dot} منفی بودن لازم نيست \dash{} اين کل نکته معرفی اپسيلون است.

خارج از لايه مرزی، \lr{V-dot} برای ۱۰۰ درصد نمونه‌ها منفی است. تضمين پايداری به ناحيه خارج از لايه مرزی اعمال می‌شود. داخل لايه مرزی، سيستم کران‌دار می‌ماند \dash{} واگرا نمی‌شود \dash{} اما همگرايی توسط شيب خطی تابع اشباع به جای جمله سوئيچينگ سخت اداره می‌شود. يک سيستم با يک لايه مرزی ۳ سانتی‌متری دور تعادل همچنان عملاً پايدار است.

%% ============================================================
\section{سوال پنجم: نتايج مونت‌کارلو چقدر تکرارپذير هستند؟}

\textbf{پاسخ:} روش‌شناسی کاملاً مستند است: ۱۰۰ اجرا به ازای هر کنترل‌گر، توزيع‌های يکنواخت با محدوده‌های بيان‌شده برای همه شش مؤلفه شرط اوليه، مدل غيرخطی کامل، مدت ۱۰ ثانيه، معيار موفقيت تعريف‌شده به عنوان هر دو زاويه در باند ۲ درصد برای يک ثانيه پيوسته.

يک شکاف تکرارپذيری وجود دارد و به صراحت مستند شده: دانه تصادفی استفاده‌شده برای توليد شرايط اوليه ثبت نشد. نتايج آماری \dash{} ميانگين‌ها، انحراف معيارها، فاصله‌های اطمينان \dash{} معتبر و تکرارپذير هستند به اين معنا که نمونه‌های تکراری از همان توزيع‌ها نتايج را در فاصله‌های اطمينان بيان‌شده توليد می‌کنند. اما دقيقاً همان ۱۰۰ بردار شرط اوليه بدون دانه قابل بازتوليد نيست.

آيتم اقدام ثبت شده: اجراهای معيار آتی دانه تصادفی را در شروع هر مطالعه مونت‌کارلو ثبت می‌کنند. نتيجه‌گيری‌های آماری \dash{} \lr{STA} سريع‌تر از کلاسيک، \lr{STA} لرزش کمتر، \lr{d} کوهن ۲.۱۴ \dash{} در برابر اين شکاف مقاوم هستند.

%% ============================================================
\section{سوال ششم: چرا ۳.۷ درصد را به جای ۶۶.۵ درصد باور کنيم؟}

\textbf{پاسخ:} معيار اصلی \lr{RMS} را روی پنجره شبيه‌سازی کامل محاسبه کرد. برای يک اجرای پايدارسازی، ۲ تا ۳ ثانيه اول فاز دسترسی هستند که در آن سيگنال کنترل صرف نظر از اپسيلون بزرگ و سريع‌تغيير است. \lr{RMS} تحت سلطه اين گذرا بود، نه لرزش حالت ماندگار که لايه مرزی برای آن طراحی شده.

معيار اصلاح‌شده \lr{FFT} را به بخش حالت ماندگار سيگنال اعمال می‌کند و انرژی بالای ۲۰ هرتز را اندازه می‌گيرد. اين مستقيماً لرزش را جدا از ديناميک‌های دسترسی اندازه‌گيری می‌کند.

استدلال فيزيکی: کاهش لرزش ۶۶.۵ درصدی از تغيير يک پارامتر اسکالر با بهره سوئيچينگ \lr{K} مساوی ۲.۲۳ و کران اغتشاش غيرصفر از نظر فيزيکی نامحتمل است. عدد ۳.۷ درصد با آنچه مکانيزم لايه مرزی بايد مشارکت کند سازگار است.

اصلاح از طريق حسابرسی داخلی شناسايی شد. توسط شخص ثالث مستقل تأييد نشده \dash{} محدوديتی که صادقانه اعتراف می‌شود.

%% ============================================================
\section{سوال هفتم: قوی‌ترين و ضعيف‌ترين نقاط شما چيست، و قبل از ارسال چه باقی می‌ماند؟}

اين سوال صداقت را پاداش می‌دهد. عدم قطعيت يا مبهم‌گويی اينجا بدتر از اعتراف صادقانه به شکاف‌هاست.

\textbf{برای نقاط قوت:} مدل سيستم صريح با همه پارامترهای فيزيکی، مشتق‌شده از مکانيک لاگرانژی. چهار فرمول‌بندی کنترل‌گر با اشتقاق‌های کامل. اثبات‌های لياپانوف کامل برای همه شش کنترل‌گر با تصديق عددی. يک معيار مونت‌کارلو ۴۰۰ شبيه‌سازی با فاصله‌های اطمينان \lr{bootstrap}، مقايسه‌های ولچ از پيش برنامه‌ريزی‌شده، و اندازه اثر \lr{d} کوهن. تنظيم \lr{PSO} استحکام با ۱۵ سناريو که تخريب لرزش را از ۵۰ برابر به ۷.۵ برابر کاهش داد. مستندسازی شفاف همه محدوديت‌ها، موارد باز، و داستان‌های اصلاح.

\textbf{برای نقاط ضعف:} سه مورد خاص. اول، بهره‌های استحکام \lr{STA} قيد نظری را نقض می‌کنند و نياز به بهينه‌سازی مجدد محدودشده يا اثبات به‌روزشده دارند. دوم، آناليز استحکام سطوح اغتشاش درشت استفاده می‌کند و يک مطالعه حساسيت کامل شناسايی‌کننده ترکيبات پارامتری خطرناک‌ترين فاقد دارد. سوم، جدول مقايسه ادبيات هنوز کامل نشده \dash{} کار در برابر کنترل‌گرهای \lr{DIP} منتشرشده معيارسنجی نشده.

\textbf{برای آنچه باقی می‌ماند:} شش آيتم با تاريخ هدف. گسترش کتاب‌نامه با ۱۵ يا بيشتر مرجع. آناليز شکاف شبيه‌سازی به سخت‌افزار. زمان‌های ساعت ديوار \lr{PSO}. شرايط اوليه \lr{QW-2} اضافه‌شده به متن اصلی. يکپارچه‌سازی ارجاع شکل برای چهار شکل کامل‌شده. بخش کار آينده اختصاصی که همه موارد باز را جمع می‌کند. همه اينها تا بازبينی مشاور ۲۴ ژانويه هدف‌گذاری شده‌اند.

%% ============================================================

\begin{mdframed}[
  backgroundcolor=audiotan,
  linecolor=audionote!70,
  linewidth=1pt,
  innerleftmargin=12pt,
  innerrightmargin=12pt,
  innertopmargin=8pt,
  innerbottommargin=8pt
]
\textbf{\color{audionote}[يادداشت صوتی: خلاصه ۶۰ ثانيه‌ای]}
اگر مشاور خلاصه يک دقيقه‌ای بخواهد: گزارش يک چارچوب شبيه‌سازی پايتون را مستند می‌کند که چهار کنترل‌گر مُد لغزشی برای يک آونگ معکوس مضاعف پياده‌سازی می‌کند، با \lr{PSO} بهينه‌سازی شده و از طريق معيار مونت‌کارلو تأييد شده. معادلات حرکت لاگرانژی صريح با اتصال کامل ارائه شده. همه چهار کنترل‌گر اثبات‌های پايداری لياپانوف کامل دارند. \lr{PSO} بهره‌های اسمی با شرايط نظری معتبر و بهره‌های استحکام با پوشش سناريوی بهبودیافته يافت. معيارهای مونت‌کارلو نشان می‌دهند \lr{STA} سريع‌ترين يکنواختی، کمترين لرزش، و بالاترين رد اغتشاش را دارد، با نتايج آماری معنادار و اندازه‌های اثر بزرگ. سه شکاف شناخته‌شده باقی می‌ماند: نقض بهره استحکام \lr{STA}، مقايسه ادبيات، و آناليز شبيه‌سازی به سخت‌افزار. اينها شناسايی شده‌اند با برنامه‌های تکميل و تاريخ‌های هدف. مقاله تحقيقاتی آماده ارسال است و پايان‌نامه در ۹۰ از ۱۰۰ به سوی ارسال نهايی است.
\end{mdframed}

\vspace{0.6cm}

%% ============================================================
\vspace{1.5cm}
\begin{center}
{\color{seriesblue}\rule{0.45\linewidth}{0.4pt}}\\[8pt]
{\small\color{sectiongray}
منابع گزارش: گزارش کامل مشاور \dash{} بخش‌های ۱ تا ۷ و يادداشت‌های سازگاری متقاطع
}
\end{center}

\end{document}
